\documentclass[12pt,a4paper]{article}
\usepackage[T1]{fontenc}
\usepackage[utf8]{inputenc}
\usepackage[margin=25mm]{geometry}
\usepackage{graphicx,booktabs,amsmath,array,microtype,xcolor}
\usepackage{tikz}
\usepackage[labelfont=bf,font=small]{caption}
\usepackage[hidelinks]{hyperref}
\usepackage{float,fancyhdr}
\setlength{\parindent}{0pt}
\setlength{\parskip}{5pt plus 1pt}
\fancyhf{}\fancyfoot[C]{\thepage}\renewcommand{\headrulewidth}{0pt}\pagestyle{fancy}
\newcommand{\resultsroot}{../output/final-results/}
\newcommand{\panel}[3]{\begin{minipage}[t]{#1}\centering\includegraphics[width=\linewidth]{\resultsroot#2}\\[2pt]{\scriptsize #3}\end{minipage}}
\newcommand{\smallpanel}[3]{\begin{minipage}[t]{#1}\centering\includegraphics[width=\linewidth,height=22mm,keepaspectratio]{\resultsroot#2}\\[2pt]{\scriptsize #3}\end{minipage}}
\begin{document}
\begin{center}
{\Large\bfseries Pool-ball localisation and recognition}\\[4pt]
Computer Vision (CSU44053)\\
Manon Yiqin Galian Ng (23375206)
\end{center}
\section*{Introduction}
The aim was to locate pool balls in photographs and identify their type and number.
I used \texttt{AllBalls.jpeg} to select detection parameters and build reference colour
models, then repeatedly evaluated and refined the approach on Tables1--6.

\section{Theory and application}
\subsection{Table masking and ball detection}
The pipeline detects circles, extracts ball crops, estimates type and colour, assigns
identities jointly, then transforms the results into an overhead view. Detection and
colour analysis use the original photograph, since rectifying it first could distort
circular ball outlines.

The supplied table corners, clockwise from top-left, are $(337,35)$, $(1570,21)$,
$(1888,451)$ and $(38,472)$. I used their quadrilateral to mask the background,
converted the image to greyscale and applied a $5\times5$ Gaussian filter with
$\sigma=1.5$ to reduce small intensity variations.

I used a gradient-based Circular Hough Transform to estimate centres and radii.
After trying parameters manually, I searched 480 combinations, keeping radii within
15--30 pixels. Configurations were ranked by matched balls, then false positives,
then mean centre error. Recall came first because a missed ball cannot be recognised
later. The selected accumulator resolution ratio was 1.5, minimum centre separation
18 pixels, edge threshold 75 and circle-vote threshold 35.

Pockets can also produce circular detections. I transformed candidate centres into
table coordinates and rejected those within 25 units of six approximate pocket
centres. Each remaining ball was cropped from the original colour image and given
a circular mask for recognition.
\begin{figure}[H]\centering
\panel{.48\linewidth}{mainDetectionResults/AllBalls/masked.png}{Table mask}\hfill
\panel{.48\linewidth}{mainDetectionResults/AllBalls/blurred.png}{Greyscale and smoothing}\\[4pt]
\panel{.48\linewidth}{mainDetectionResults/AllBalls/raw_circles.png}{Hough circles}\hfill
\panel{.48\linewidth}{mainDetectionResults/AllBalls/filtered_circles.png}{Pocket filtering}
\caption{Detection stages on AllBalls. Pocket filtering removes the extra circle.}
\end{figure}

\newpage
\subsection{Ball type and body-colour histograms}
I clustered the pixels inside each circular crop in BGR colour space using $K=2$.
The idea was that a striped ball's white and coloured regions would often form
separate groups. I used three initialisations with a fixed random seed, stopping
after 50 iterations or when cluster centres changed by less than 0.2. Each cluster's
mean colour was converted to HSV: $V\ge170$ and $S\le165$ indicated light pixels;
a non-light cluster with $V\le90$ was dark. These thresholds were chosen manually.
Cluster sizes then estimated the light fraction $w$.

For colour recognition, I built eight reference HSV histograms from labelled
AllBalls examples, sharing a model between solid and striped partners and excluding
the cue. A $3\times3$ elliptical erosion reduced boundary pixels, and bright,
low-saturation pixels were removed to suppress white patches and reflections.
Crops with fewer than 12 usable body pixels supplied no numbered-colour evidence.
\begin{figure}[H]\centering
\begin{minipage}[t]{.24\linewidth}\centering
% Original-image crop centred on Table3 candidate 8; not a new application export.
\begin{tikzpicture}\clip (0,0) circle (11mm);
\node[inner sep=0] at (0,0) {\includegraphics[width=22mm,viewport=1392 444 1431 483,clip]{\resultsroot mainDetectionResults/Table3/masked.png}};
\end{tikzpicture}\\[2pt]{\scriptsize Circular colour crop}
\end{minipage}\hfill
\smallpanel{.24\linewidth}{mainTypeRecognitionResults/Table3/ball_8_kmeans.png}{Two colour clusters}\hfill
\smallpanel{.24\linewidth}{mainColourRecognitionResults/Table3/ball_8_body.png}{Retained body pixels}\hfill
\smallpanel{.24\linewidth}{ExperimentHistogramBackprojectionResults/Table3/ball_8_bestColour.png}{Reference backprojection}
\caption{Recognition stages for Table3 candidate 8. Backprojection visualises colour agreement; recognition uses direct histogram comparison.}
\end{figure}

Each histogram has $18\times8\times8=1152$ bins. Hue spans $[0,180)$, and
saturation and value span $[0,256)$, so bins cover 10 hue units and 32 saturation
or value levels. Finer hue grouping helps distinguish body colours, while broader
saturation and value groups allow lighting variation. This is a design compromise,
not a claim that these bin counts are optimal.

For nearly grey pixels (BGR maximum minus minimum below 15), I set hue to zero
because it is unstable at low chroma. A $[1,2,1]/4$ filter smooths neighbouring bins,
wrapping hue and repeating the other endpoints. Histograms sum to one, with solid
and striped references weighted equally. Candidate histogram $p$ is compared with
each reference $q$ using one minus Bhattacharyya distance [1]:
\[
s(p,q)=1-\sqrt{1-\sum_b\sqrt{p_bq_b}}.
\]
This produces eight colour similarities rather than one fixed colour label.

\newpage
\subsection{Type compatibility and identity assignment}
Hard thresholds can give different types to balls with nearly identical light
coverage. I therefore used soft compatibility scores with expected light-fraction
intervals $[0,0.30]$ for solids, $[0.30,0.85]$ for stripes and $[0.85,1]$ for cue.
For type $t$ and observed light fraction $w$:
\[
g_t(w)=\exp\left(-\frac{d(w,I_t)^2}{2(0.10)^2}\right).
\]
Here $d$ is zero inside interval $I_t$ and the distance to its nearest endpoint
outside it. Compatibility is one inside the interval and decreases smoothly outside;
the width 0.10 controls the decrease. This is a heuristic score, not a probability.
Numbered-identity scores multiply colour similarity by type compatibility. The cue
score is $w g_{\mathrm{cue}}(w)$.

Each ball identity can appear at most once. I arranged candidate-to-identity scores
in a matrix and used Hungarian assignment [2] to maximise their total while enforcing
uniqueness. Extra dummy columns allow several candidates to receive \texttt{Unknown},
mainly for surplus candidates or those without usable evidence. This is not a
calibrated confidence-based rejection rule.

Sometimes a detected circle enclosed a white interior patch rather than the whole
ball. I expanded the recognition radius of initially white candidates by one pixel
and recomputed their scores, keeping the original detection centre and radius.
If several candidates were initially white, I reserved cue for the largest original
Hough radius. This uses size as a clue to distinguish a whole cue ball from
a smaller white patch.

\begin{figure}[H]\centering
\begin{minipage}[t]{.48\linewidth}\centering
\includegraphics[height=16mm,viewport=244 309 292 357,clip]{\resultsroot mainDetectionResults/Table1/masked.png}\\[2pt]
{\scriptsize Genuine cue: original radius 24.15 pixels}
\end{minipage}\hfill
\begin{minipage}[t]{.48\linewidth}\centering
\includegraphics[height=16mm,viewport=1330 280 1366 316,clip]{\resultsroot mainDetectionResults/Table1/masked.png}\\[2pt]
{\scriptsize Competing patch: original radius 17.55 pixels}
\end{minipage}
\caption{Initially white candidates on Table1, shown at equal display size. The smaller circle captures an interior region of a striped ball.}
\end{figure}

\subsection{Coordinate transformation and visualisation}
A homography maps the four table corners to $(0,0)$, $(1433,0)$, $(1433,682)$ and
$(0,682)$ in a $1434\times683$ rectangle:
\[
\lambda(x_t,y_t,1)^T=H(x,y,1)^T.
\]
A table unit is approximately one millimetre at the supplied scale. Stored centres
are transformed for matching and output. The overhead diagram uses fixed ball radii
and colours/bands representing predicted identities, rather than observed orientation.
\begin{figure}[H]\centering
\panel{.32\linewidth}{mainFinalResults/Table2/original_numbered.png}{Original-image labels}\hfill
\panel{.32\linewidth}{mainFinalResults/Table2/table_numbered.png}{Rectified labels}\hfill
\panel{.32\linewidth}{mainFinalResults/Table2/table_diagram.png}{Overhead diagram}
\caption{Transformation and presentation on Table2. All views retain the extra false detection.}
\end{figure}

\newpage
\subsection{Alternative approaches investigated}
\textbf{Fixed HLS thresholds.} My first type estimator separated pixels into white,
dark and coloured groups using HLS lightness and saturation. Their proportions
supplied cue/solid/stripe rules. These were easy to interpret but sensitive to
highlights and shadows, motivating K-means. The initial K-means rules
labelled at least 85\% light coverage as cue, at least 70\% non-light coverage as
solid, and other mixed cases as striped. I also tested $K=3$--6.

\textbf{Connected-component analysis (CCA).} Colour proportions do not describe
spatial arrangement. I therefore examined component counts, largest-region fractions
and boundary light coverage. Larger white regions might distinguish stripes from
solid balls' number patches, but highlights, orientation and thresholded shadows
can change connectivity. Initial colour recognition used cluster-mean hue
thresholds, with a separate low-value rule for black. Similar colours and shifting
cluster means motivated the histogram descriptor.
\begin{figure}[H]\centering
\panel{.37\linewidth}{ExperimentManualHLSResults/Table1/predicted_types.png}{HLS type predictions}\hfill
\panel{.37\linewidth}{ExperimentKMeansK2Results/Table1/predicted_types.png}{Initial K=2 predictions}\hfill
\smallpanel{.19\linewidth}{ExperimentCCAResults/Table1/ball_7_light.png}{CCA light region}
\caption{Alternative type evidence on Table1. A light region alone does not establish that a ball is striped or cue.}
\end{figure}

\textbf{Segmentation-assisted detection.} Touching balls interrupt each other's
outlines. I tested foreground segmentation and distance-transform/watershed proposals
near baseline detections, using area, shape, support, duplicate and pocket checks.
A separate combined method estimated local cloth colour and routed foreground
clusters to more permissive local Hough searches, retaining baseline circles outside
these regions. Edge, arc and foreground checks filtered proposals; a subsequent
marker-controlled watershed rejected circles with implausible supported basin areas.
These checks used image evidence rather than ground-truth labels.
\begin{figure}[H]\centering
\panel{.32\linewidth}{ExperimentWatershedResults/Table1/foreground.png}{Foreground}\hfill
\panel{.32\linewidth}{ExperimentWatershedResults/Table1/distance.png}{Distance transform}\hfill
\panel{.32\linewidth}{ExperimentWatershedResults/Table1/basins.png}{Watershed basins}\\[4pt]
\panel{.48\linewidth}{ExperimentGuidedHoughWatershedResults/Table3/clusters.png}{Combined method: local search regions}\hfill
\panel{.48\linewidth}{ExperimentGuidedHoughWatershedResults/Table3/basins.png}{Combined method: watershed basins}
\caption{Segmentation stages. Separating a region visually does not guarantee an accepted ball detection.}
\end{figure}

\newpage
\section{Results and analysis}
\subsection{Evaluation methodology}
Tables1--6 contain 90 annotated balls. Predictions are matched one-to-one to annotated
centres within 15 rectified table units. Matching prioritises the number of valid
pairs, using proximity to resolve alternatives; it need not minimise total distance
globally. Matched detections are true positives (TP), unmatched predictions false
positives (FP), and unmatched annotations false negatives (FN). Pooled counts give:
\[
P=\frac{TP}{TP+FP},\qquad R=\frac{TP}{TP+FN},\qquad F_1=\frac{2PR}{P+R}.
\]
Mean localisation error is $\frac{1}{TP}\sum_i\|\hat{\mathbf{x}}_i-\mathbf{x}_i\|_2$
for matched centres in table coordinates. Type accuracy measures cue/solid/striped
correctness; identity accuracy requires the exact number or cue. If $C$ identities
are correct, matched identity accuracy is $C/TP$ and end-to-end correct-identity
recall is $C/90$. An Unknown label on a genuine ball is an identity error.
Extra detections affect precision, but are not directly penalised by $C/90$.

\subsection{Final performance on the supplied images}
\begin{table}[H]\centering\small
\caption{Final detection and recognition results. Mean error is in table units.}
\begin{tabular}{lrrrrrr}\toprule
Image & GT & TP & FP & FN & Mean error & Correct identities\\\midrule
Table1 &16&14&1&2&7.68&11/14\\
Table2 &16&16&1&0&3.58&16/16\\
Table3 &16&15&0&1&3.12&15/15\\
Table4 &15&14&0&1&3.89&14/14\\
Table5 &14&14&0&0&5.83&13/14\\
Table6 &13&13&0&0&5.76&11/13\\\midrule
Total &90&86&2&4&4.91&80/86\\\bottomrule
\end{tabular}
\end{table}
The detector achieves 97.73\% precision, 95.56\% recall and 96.63\% F1. Identity
accuracy is 93.02\% on matched balls, while end-to-end correct-identity recall is
88.89\%. The latter reflects the complete task because missed balls also
count against it. All six genuine cues are identified correctly. The pooled mean
centre error is 4.91 units, weighting images by their number of matches.

Table2 recognises all genuine balls but retains a false detection. Tables3 and 4
recognise every matched ball, yet each misses ball 10. Tables5 and 6 detect every
ball but still make identity mistakes. Detection and recognition therefore remain
separate sources of failure; high accuracy on matched balls does not imply that
the whole scene is correct.

On AllBalls, the selected search configuration gave 16 TP, 1 FP and no FN
(94.12\% precision, 100\% recall, 96.97\% F1; mean error 2.76).
Pocket filtering removed the false positive, and final recognition was 16/16.
These are training-image results, since AllBalls supplied parameter selection and
reference colours; they are excluded from the development totals.

\newpage
\subsection{What the experiments changed}
\begin{table}[H]\centering\small
\caption{Detection alternatives on Tables1--6, using pooled counts.}
\begin{tabular}{lrrrr}\toprule
Method & TP & FP & FN & F1 (\%)\\\midrule
Selected global Hough + filtering &86&2&4&96.63\\
Watershed augmentation &86&2&4&96.63\\
Cluster-local Hough &86&16&4&89.58\\
Cluster-local Hough + watershed &83&11&7&90.22\\\bottomrule
\end{tabular}
\end{table}
Watershed augmentation accepted no detection changes. Local Hough retained the
same total TP count but produced 14 more false positives than the selected detector.
Its watershed check removed five false positives while losing three genuine matches.
Although this improved F1 slightly over local Hough alone, both remained worse than
the simpler global detector. A local recovery therefore did not establish an overall
improvement.
\begin{figure}[H]\centering
\panel{.32\linewidth}{mainDetectionResults/Table3/filtered_circles.png}{Baseline: 15 TP, 0 FP, 1 FN}\hfill
\panel{.32\linewidth}{ExperimentGuidedHoughPreWatershedResults/Table3/circles.png}{Local Hough: 16 TP, 4 FP, 0 FN}\hfill
\panel{.32\linewidth}{ExperimentGuidedHoughWatershedResults/Table3/circles.png}{Watershed: 16 TP, 3 FP, 0 FN}
\caption{Table3 illustrates a recovered ball accompanied by false detections.}
\end{figure}
\begin{table}[H]\centering\small
\caption{Recognition comparisons. Type and identity counts measure different tasks.}
\begin{tabular}{llr}\toprule
Task & Configuration & Correct /86\\\midrule
Type & HLS thresholds &43\\
Type & K=2 &76\\
Type & K=3, 4 or 5 (each) &73\\
Type & K=6 &74\\
Type & K=2 + CCA &73\\\midrule
Identity & Hue rules &68\\
Identity & HSV histograms, hard types &72\\
Identity & Independent soft scores &73\\
Identity & Add unique assignment &75\\
Identity & Add selective +1 crop expansion &79\\
Identity & Add largest-radius cue reservation &80\\\bottomrule
\end{tabular}
\end{table}
K=2 outperformed the tested larger cluster counts under the same type rules. Extra
clusters can represent shading rather than useful material regions. CCA also reduced
type correctness, so I retained the simpler estimator. The original controlled
recognition comparisons held detections fixed: histograms added four correct
identities, soft scoring one, assignment two, selective expansion four and cue
reservation one. The total gain cannot be credited to histograms or assignment alone.

\newpage
\subsection{Successes and remaining failures}
\begin{figure}[H]\centering
\panel{.32\linewidth}{mainFinalResults/Table1/table_numbered.png}{Table1: 11/14 identities; two misses}\hfill
\panel{.32\linewidth}{mainFinalResults/Table2/table_numbered.png}{Table2: 16/16; one false detection}\hfill
\panel{.32\linewidth}{mainFinalResults/Table3/table_numbered.png}{Table3: 15/15; one miss}\\[5pt]
\panel{.32\linewidth}{mainFinalResults/Table4/table_numbered.png}{Table4: 14/14; one miss}\hfill
\panel{.32\linewidth}{mainFinalResults/Table5/table_numbered.png}{Table5: 13/14; no misses}\hfill
\panel{.32\linewidth}{mainFinalResults/Table6/table_numbered.png}{Table6: 11/13; no misses}
\caption{Final labelled results for all six development images. Fractions count correct identities among matched detections; labels include false detections and Unknown.}
\end{figure}

\textbf{Crowding and localisation.} Table1 misses balls 6 and 7 in the crowded group,
assigns a false circle to 6, leaves ball 14 Unknown and swaps identities 5 and 9.
Interrupted outlines make circle evidence weaker in crowded regions. The swapped
matches have centre errors of 14.63 and 14.81 units, close to the 15-unit tolerance,
so their recognition errors coincide with poor localisation. A displaced crop can
include neighbouring colour or exclude part of a stripe, giving unreliable evidence
even when a centre still counts as matched. Across all matches, median error is
3.75 units, maximum error is 14.81, and seven errors exceed 10 units.

\textbf{Colour and type confusion.} Tables5 and 6 label ball 10 as 4, changing both
blue to purple and stripe to solid. Table6 also labels 3 as 11, preserving red but
changing solid to stripe. The latter shows that correct body colour is not enough
for correct identity. Stripe orientation can hide much of the white region, while
reflections can make a solid appear lighter. These are plausible explanations for
the confusions, rather than verified causes for each individual error.

\textbf{Cue correction.} Selective expansion improved the 75/86 control to 79/86
without losses, and cue reservation added the final correction. The row-adjusted
radius rule made the same choices. Expanding every ball had worsened a separate
baseline, so the benefit is specific to initially white crops. Assignment can also
propagate a score change to other candidates.

\begin{figure}[H]\centering
\begin{minipage}[t]{.32\linewidth}\centering
\includegraphics[height=24mm,viewport=1010 203 1275 488,clip]{\resultsroot mainFinalResults/Table1/table_numbered.png}\\[2pt]
{\scriptsize Table1: crowding and false circle 6}
\end{minipage}\hfill
\begin{minipage}[t]{.32\linewidth}\centering
\includegraphics[height=24mm,viewport=1225 473 1320 563,clip]{\resultsroot mainFinalResults/Table6/table_numbered.png}\\[2pt]
{\scriptsize Table6: ball 10 labelled 4}
\end{minipage}\hfill
\begin{minipage}[t]{.32\linewidth}\centering
\includegraphics[height=24mm,viewport=1165 198 1245 283,clip]{\resultsroot mainFinalResults/Table6/table_numbered.png}\\[2pt]
{\scriptsize Table6: ball 3 labelled 11}
\end{minipage}
\caption{Enlarged details of the final outputs make the remaining failures visible.}
\end{figure}

\newpage
\subsection{Robustness, appropriateness and improvements}
The selected combination balances accuracy and complexity on the supplied images.
Circular detection suits approximately round, visible balls, histogram distributions
are less dependent on one representative colour, and uniqueness uses a real constraint
of the task. The rejected alternatives did not improve their tested controls; this
does not mean those methods are ineffective in other settings.

\textbf{Limits of the evidence.} Tables1--6 share a camera and broadly similar
conditions, with some repeated arrangements. Repeated inspection and method selection
introduce selection bias, even though histogram references came only from AllBalls.
The reported scores describe development performance rather than performance on an
untouched test set. One additional correct identity changes matched accuracy by
about 1.16 percentage points, so small gains need further testing.

\textbf{Changes in scene conditions.} Fixed image-space radii, table corners and
pocket locations depend on the camera setup. A new viewpoint or resolution would
require updated geometry and radius ranges. Occlusion can remove the circular
boundary needed by Hough detection. Different lighting, cloth colour or camera colour
response can change histogram matches, and stripe orientation changes visible white
coverage. The largest-radius cue rule can fail if a distant cue appears smaller than
a nearby white patch. Evaluating new photographs with varied viewpoints, lighting
and crowded arrangements, while freezing settings before inspection, is therefore
the most useful next step.

\textbf{Recognition improvements.} More labelled reference examples under varied
illumination and orientation could make the colour models less dependent on one
training image. Better boundary estimation or cloth rejection could reduce crop
contamination, but should be checked for both fixes and regressions. A calibrated
Unknown policy could prevent weak candidates from taking real identities during
assignment. The current best-effort assignment improves results, but its uniqueness
constraint can also spread one candidate's error to another.

\textbf{Evaluation improvements.} The 15-unit tolerance is about 35\% of the 43 mm
ball diameter. Reporting performance across several tolerances would show how much
success depends on accepting weakly localised matches; it should not be adjusted
simply to raise the score. The homography describes the cloth plane, while visible
ball centres lie above it, creating possible systematic error. Centre annotations
also cannot establish the accuracy of masks or full ball boundaries. Boundary labels
and isolated recognition comparisons would help distinguish detection geometry,
colour evidence and assignment errors.

\paragraph{References}
\small
[1] K. Dawson-Howe, CSU44053 course notes: \emph{Images, Edges, Histograms, Region
Segmentation, Learning and Evaluation}; based on \emph{A Practical Introduction to
Computer Vision with OpenCV}, Wiley (2014).\par
[2] H. W. Kuhn, ``The Hungarian method for the assignment problem,''
\emph{Naval Research Logistics Quarterly}, 2, 83--97 (1955).
\end{document}
